% compendium/digital-assets-prediction-markets/crypto-venues-perpetuals.tex
% Entries migrated on 2026-10-15 from the legacy crypto-market-microstructure file (slice 24), owned by
% crypto-perpetuals (a routing exception from the file's map candidate crypto-mm-microstructure);
% dispositions and source checks in notes/migration.md, the blocks verbatim in
% notes/artefacts/crypto-mm-microstructure/derivations.tex.

\section{Crypto venues, perpetuals and trading risk}

\begin{qfdefinition}[Crypto venue and instrument structure: fragmented executable prices, reference prices and best execution, and spot, dated futures, perpetuals (linear, inverse, quanto), options and pools]\label{def:crypto-venue-instrument-structure}
\qffield{Statement} Crypto assets trade on many independently owned venues. Makarov and Schoar, using tick data from 34 exchanges in 19 countries, find that individual exchanges mostly run centralised order books like equity markets, but that, unlike regulated equity markets, the crypto market ``lacks any provisions to ensure that investors receive the best price'', contrasting the U.S. NBBO requirement \citep[discussion-paper version of December 2018, pp.~1--2 and footnote~1]{s6-makarov}. Kept apart here:
\begin{itemize}
\item \emph{executable prices}: the bids and offers on each venue, in its own quote currency, custody and fee terms;
\item \emph{reference prices}: indices a venue or benchmark administrator computes from several venues, against which derivatives are marked or settled (Deribit marks positions with a Mark Price and computes funding against a Deribit Index, \hyperref[mech:perpetual-funding]{perpetual funding});
\item \emph{fair-value models}: a trader's own estimate of value, which need not equal either;
\item \emph{legal best-execution obligations}: duties under a jurisdiction's rules (\hyperref[def:best-execution]{best execution}), which Makarov and Schoar found absent across crypto venues as a market-wide price-protection rule (p.~1).
\end{itemize}

\emph{Instruments.} \emph{Spot} trades exchange one asset for another; on Deribit spot trades are brokerage services settled and delivered immediately \citep[Rulebook v2.0, Rules 1.2(a) and 11.3(a)]{gapclose-deribit-rules}. \emph{Dated futures} expire, and convergence to spot at expiry limits the basis \citep[arXiv v2, p.~2]{roadmap-perpetual-pricing}. \emph{Perpetual futures} have no expiry and are anchored by periodic funding payments (\hyperref[mech:perpetual-funding]{perpetual funding}). Ackerer, Hugonnier and Jermann distinguish three contract types (pp.~2--3):
\begin{itemize}
\item \emph{linear}: contract size in units of the base asset (1 BTC), margined and settled in the quote currency (USD, or in practice a stablecoin such as USDT, p.~7);
\item \emph{inverse}: the base asset itself (BTC) is the margin and settlement asset, and the contract size is in the quote currency (for example 10{,}000 USD);
\item \emph{quanto}: a third currency (for example ETH) is used for margining and settlement.
\end{itemize}
\emph{Options}: Deribit lists European-style cash-settled options and linear options; its European options are priced in BTC or ETH and cannot be exercised before expiry (Rulebook, Appendix~A, 2(d)). \emph{Listing snapshot, 17 October 2026} (Deribit public API, one venue on one day; open interest as the API reports it, in units of the underlying, the contract size being one coin for both BTC listings; units not confirmed from the API documentation). BTC and ETH options are listed as inverse (``reversed'', premium and settlement in the coin): 938 BTC instruments over 12 expiries, a median of 42 strikes per expiry. USDC-settled linear options are listed on BTC, ETH, SOL, XRP, AVAX, HYPE and TRX. Inverse BTC options carry 390{,}301 of open interest against 2{,}781 for USDC-settled BTC options (0.7\%). Open interest is concentrated: the top 10\% of inverse BTC instruments hold 67\% of it, 202 of 938 hold none, and two expiries (30~October and 25~December 2026) hold 72\% \citep[\texttt{get\_instruments} and \texttt{get\_book\_summary\_by\_currency}, 17 October 2026]{deribit-api-2026}; script \nolinkurl{batch60-deribit-options-snapshot.py} in \nolinkurl{notes/artefacts/crypto-perpetuals/}. \emph{On-chain pools} (automated market makers) were not checked in this slice.

\qffield{Conditions} The fragmentation evidence is historical. Over December 2017 to February 2018 the daily average price ratio between the United States and Korea exceeded 15 percent and reached 40 percent on several days; Japan--U.S. differences were about 10 percent and U.S.--Europe about 3 percent, while differences within a country typically did not exceed one percent on average (pp.~1--2). The authors estimate at least \$2 billion of potential arbitrage profit over that period from executed trades (p.~2, footnote~2). Spreads co-moved across countries and opened during rapid appreciation; between cryptocurrencies on the same exchanges spreads were an order of magnitude smaller, and pairwise capital-control tightness was positively related to the correlation of spreads, which the authors read as evidence that capital controls contribute (pp.~2--3). These are 2017--2018 findings; current spreads were not measured.

\qffield{Practice} Interpretation: a crypto price is a price on a venue, in a quote currency, with that venue's custody, collateral and fees. Quote-currency differences (USD, USDT, USDC, KRW) and settlement currency (linear against inverse) change the units of P\&L and margin. Cross-references: fragmentation without a consolidated tape in spot FX is in \hyperref[def:fx-market-structure]{spot FX market structure}, routing across venues in \hyperref[def:smart-order-routing]{smart order routing}, limits to arbitrage in \hyperref[res:arbitrage-risks]{arbitrage risks}.

\qffield{Pitfalls} The legacy principle said ``there is no universal consolidated tape, so fair value and best execution are venue-dependent''. The sources read support the absence of a market-wide best-price rule in 2018 (Makarov and Schoar) but say nothing on consolidated tapes; ``fair value is venue-dependent'' conflates venue prices with fair value. Its taxonomy's ``perpetual futures: leveraged swaps kept near spot through funding'' overstates the anchoring (\hyperref[mech:perpetual-funding]{perpetual funding}); ``options: volatility exposure concentrated in fewer venues and strikes'' is supported only for strikes and expiries, by the one-venue snapshot above; no other venue's volumes were read, so ``fewer venues'' is unchecked. ``On-chain pools: automated market makers with inventory curves and gas costs'' was not checked here (see \hyperref[def:amm-pool-mechanics]{AMM pool mechanics}). Deribit's terms are one venue's.

\qffield{See also} \hyperref[mech:perpetual-funding]{perpetual funding}; \hyperref[def:crypto-trading-risk-layers]{crypto trading risk layers}; \hyperref[def:best-execution]{best execution}; \hyperref[def:fx-market-structure]{spot FX market structure}.

\qffield{Further reading} The legacy blocks it replaces, verbatim, are in \nolinkurl{derivations.tex} in the folder \nolinkurl{notes/artefacts/crypto-mm-microstructure/}.

\qffield{Index} crypto market fragmentation; cross-exchange arbitrage; Kimchi premium; capital controls; reference price; index price; mark price; best execution in crypto; spot crypto; dated crypto futures; perpetual futures; linear perpetual; inverse perpetual; quanto perpetual; crypto options; automated market maker; quote currency; settlement currency

\qffield{Status} Verification: partly checked 2026-10-15 (\texttt{s6-makarov}, LSE FMG Discussion Paper 782 of December 2018, abstract and pp.~1--3, text layer; \texttt{roadmap-perpetual-pricing}, arXiv v2 of 3 September 2024, pp.~2--3 and~7; \texttt{gapclose-deribit-rules}, Rulebook version 2.0 applicable 13 August 2026, retrieved 15 October 2026, Rules 1.2 and~11.3 and Appendix~A; third batch of 2026-10-17: Deribit public API instrument and book-summary responses of 17 October 2026). Stability: draft. Currency: fragmentation evidence 2017--2018; Deribit terms as of October 2026.
\end{qfdefinition}

\begin{qfmechanism}[Perpetual funding: premium and interest terms, sign convention, accrual and settlement, damper and caps (Deribit as one venue), and why convergence is not guaranteed]\label{mech:perpetual-funding}
\qffield{Statement} A perpetual future has no expiry, so it cannot rely on convergence at maturity; instead, longs and shorts exchange periodic \emph{funding} payments. Ackerer, Hugonnier and Jermann describe the payment as generally the sum of a \emph{premium} part $\kappa_t(f_t-x_t)$, paid by longs when the perpetual price $f_t$ is above spot $x_t$, and an \emph{interest} part $\iota_t x_t$ set by the exchange to reflect the interest-rate differential between the two currencies; for a linear contract the long's cash flow over one period is
\[ (f_{t+1}-f_t)-\kappa_t(f_t-x_t)-\iota_t x_t \]
in the quote currency \citep[arXiv v2, sec.~3, pp.~7--8, eqs.~(4)--(6)]{roadmap-perpetual-pricing}. Funding is commonly computed every 8 hours (p.~2). In their frictionless no-arbitrage model the exchange can choose the interest factor so that the perpetual price equals spot for any sufficiently large premium rate (p.~4; Corollary~1, p.~12). That is a model result. He, Manela, Ross and von Wachter state that perpetuals ``are not guaranteed to converge to the spot price'', derive bounds with trading costs, and find empirical deviations larger than in traditional currency markets, co-moving across currencies and diminishing over time \citep[arXiv v7 of 17 September 2026, abstract]{roadmap-perpetual-fundamentals}.

\qffield{Conditions} \emph{One venue's rule: Deribit, Funding Specifications page edited 27 July 2026, retrieved 15 October 2026} \citep[Funding Specifications]{gapclose-deribit-rules}:
\begin{itemize}
\item \emph{sign}: when the funding rate is positive, longs pay shorts; when negative, shorts pay longs; Deribit charges no fee on funding, so payments net to zero between holders;
\item \emph{rate}: premium rate $P=(\text{Mark Price}-\text{Deribit Index})/\text{Deribit Index}$; a damper sets funding to zero for $|P|\le0.025\%$ and otherwise moves $P$ towards zero by $0.025\%$; a cap then bounds the rate, $\pm0.5\%$ for BTC, $\pm1.0\%$ for ETH and $\pm5.0\%$ for the USDC- and USDT-margined perpetuals:
\[ \text{rate}=\min\bigl(c,\ \max\bigl(-c,\ \max(0.025\%,P)+\min(-0.025\%,P)\bigr)\bigr); \]
\item \emph{units and cadence}: the rate is expressed per 8 hours but accrues continuously; the payment is rate $\times$ position size $\times$ (period / 8 hours), in the settlement currency of the contract: BTC in Deribit's BTC example, where a USD 100{,}000 position is 1 BTC, an inverse contract in Ackerer, Hugonnier and Jermann's sense (Interpretation), and USDC for a USDC-margined linear one;
\item \emph{settlement}: accrued funding shows at once in account equity (realised session P\&L) and can be used for trading, but moves to the cash balance only at the daily settlement at 08:00 UTC, after which it can be withdrawn (also Rulebook v2.0, Rule 11.4).
\end{itemize}

\qffield{Practice} Worked example (Deribit's, recomputed by the script \nolinkurl{funding_checks.py} beside the artefact): BTC perpetual mark USD 100{,}075 against index USD 100{,}000 gives $P=0.075\%$ and a rate of $0.05\%$ per 8 hours; a long of USD 100{,}000 (1 BTC) pays $0.05\%\times1\text{ BTC}\times1/480=0.000001041667$ BTC per minute and $0.0005$ BTC over 8 hours. A USDC-margined SOL perpetual at mark 100.10 against index 100 gives $0.075\%$, so a USDC 100{,}000 long pays 0.15625 USDC per minute and 75 USDC over 8 hours. The script also checks the printed max/min formula against the prose rule on a grid of premia and the other examples (damper, caps, offsetting minutes). Interpretation: funding is a carry cost or income of holding the perpetual against spot, in the contract's settlement currency; whether it signals crowding was not checked.

\qffield{Pitfalls} The legacy definition said funding is ``the periodic payment that links perpetual futures to spot'', with longs paying when it is positive, and that ``it is both a carry term and a crowding signal''. The sign convention matches Deribit's; ``links \dots\ to spot'' holds as a design intent and a model result, not a guarantee; ``crowding signal'' is not supported by the sources read. Third batch of 2026-10-17: the dated-futures evidence in \hyperref[ev:crypto-dated-futures-carry]{crypto carry} ties the basis to small traders' demand for leverage and to limited arbitrage capital, which supports a leverage-demand reading. Direct evidence on funding rates was again not read (\texttt{ev:crypto-funding-crowding}). Venue rules differ: other venues' cadence, interest term, clamps and index construction were not checked, and Deribit's own page contains slips in its ETH examples (an index ``at USD 10,000'' and ``a 1 BTC long position'' in the ETH column of Example~3; 7.5\% inside the ETH cap formula of Example~5), whose stated results are nevertheless correct.

\qffield{See also} \hyperref[def:crypto-venue-instrument-structure]{crypto venue and instrument structure}; \hyperref[def:crypto-trading-risk-layers]{crypto trading risk layers}; \hyperref[mech:futures-margin-settlement]{futures margin and settlement}; \hyperref[def:securities-lending-fee-rebate]{lending fees and rebates} (financing legs).

\qffield{Further reading} The legacy block it replaces, verbatim, is in \nolinkurl{derivations.tex} in the folder \nolinkurl{notes/artefacts/crypto-mm-microstructure/}.

\qffield{Index} perpetual futures; funding rate; funding payment; premium rate; interest rate differential; anchoring; mark price; Deribit index; funding damper; funding cap; 8-hour funding rate; continuous funding accrual; daily settlement; linear perpetual funding; inverse perpetual funding; perpetual basis

\qffield{Status} Verification: partly checked 2026-10-15 (\texttt{roadmap-perpetual-pricing}, arXiv v2 of 3 September 2024, pp.~2, 4, 7--8 and~12, text layer; \texttt{roadmap-perpetual-fundamentals}, arXiv v7 of 17 September 2026, abstract; \texttt{gapclose-deribit-rules}, Funding Specifications edited 27 July 2026 and Rulebook v2.0 Rule 11.4, retrieved 15 October 2026 through the help-centre API; formula and examples recomputed by script). Stability: draft. Currency: Deribit terms as of October 2026; model results timeless.
\end{qfmechanism}

\begin{qfdefinition}[Crypto trading risk layers: liquidation, insurance fund, socialised loss and auto-deleveraging at one venue, kept apart from custody, stablecoin and network-settlement risk]\label{def:crypto-trading-risk-layers}
\qffield{Statement} The losses a crypto trader can bear beyond market risk come in layers that are governed by different documents and fail in different ways:
\begin{itemize}
\item \emph{liquidation}: forced closing of positions when margin falls short. Under Deribit's rulebook, if a member fails to meet maintenance margin, Deribit may close positions, cancel orders, delta-hedge with futures and perpetuals, apply auto-deleveraging where applicable, or close the whole account, together the Liquidation Process, and may charge a liquidation fee \citep[Rulebook v2.0, Rules 10.1 and~10.5]{gapclose-deribit-rules};
\item \emph{insurance fund}: if an account's margin balance in a currency is negative after liquidation, Deribit transfers assets from its insurance fund, funded mainly by liquidation fees and excess collateral from liquidations, to bring it to zero; if the fund in that currency is depleted, Deribit may, in its absolute discretion, convert another currency of the fund (Rules 10.6--10.7; definition of Insurance Fund);
\item \emph{socialised loss}: Deribit may, in its absolute discretion, trigger the Socialised Loss Mechanism if the insurance fund is depleted in all currencies, or if it is depleted in the bankrupt account's currency and Deribit decides not to convert other currencies (Rule 10.8), so the mechanism can be triggered while other currencies remain in the fund, and need not be triggered when the fund is exhausted. Once it is, Deribit may deduct from all members with positive profit and loss in the same day's session, in any currency, their pro-rata share of the negative balance, halt trading, cancel orders and close positions at the mark price; members agree to this and may dispute only a manifest error in their share (Rules 10.9--10.11);
\item \emph{auto-deleveraging}: a step in the liquidation waterfall in which, once liquidity-support providers' capacity is used, eligible positions of a distressed account are transferred to a deleveraged counterparty at a preferential price and closed against that member's existing positions (definition of Auto-Deleveraging; the full waterfall is in the Exchange Membership Terms, clause~12, not read);
\item \emph{custody}: Deribit holds members' assets in its own wallets or through a third-party custodian of its choice, or acts as limited agent over a member-chosen custodian for margining, liquidation and settlement (Rule~2, guidance);
\item \emph{stablecoin, network and bridge risk}: redemption and depeg of stablecoin collateral, chain congestion, reorganisations and transfer delays. \emph{One depeg documented} (third batch of 2026-10-17). In March 2023 Silicon Valley Bank entered receivership on Friday 10 March at 11:37~a.m.\ ET. Circle announced around 10~p.m.\ that it could not withdraw \$3.3 billion of USDC reserves held there, about 8 percent of the total. Redemptions had already surged before the announcement. Primary issuance and redemption, ``constrained by working hours of the U.S. banking systems'', largely stopped over the weekend, and USDC fell to 86 cents on secondary markets. Dai's one-to-one Peg Stability Modules against USDC were drained, and stablecoins with no SVB exposure also lost their pegs. The depegs eased only after the FDIC, Treasury and Federal Reserve announced full protection of SVB depositors \citep[FEDS Notes, 17 December 2025, summary and sec.~1]{feds-2025-svb-stablecoins}. So a fully backed stablecoin used as collateral can lose its peg when its primary market is closed and its reserves are frozen, and smart-contract links can pass the loss to other coins. Chain congestion, reorganisations and bridge delays were not checked.
\end{itemize}

\qffield{Conditions} All rules above are Deribit FZE's, a virtual-asset service provider regulated by Dubai's VARA, under Rulebook version 2.0 applicable 13 August 2026 (Rule 1.1); other venues' waterfalls were not checked and may differ. Venue insolvency is a separate legal process from a venue's socialised-loss mechanism; neither insolvency outcomes nor customer-asset treatment in insolvency were checked.

\qffield{Practice} Interpretation: keep apart a trader's own liquidation (its positions closed), a loss to the venue's backstop (insurance fund), a loss charged to other members (socialised loss), a position closed against a counterparty's will (auto-deleveraging) and a loss of assets held at the venue (custody, insolvency). Cross-venue strategies add transfer and settlement paths between venues, during which capital cannot respond. Margin and default waterfalls in a regulated futures clearing house are in \hyperref[mech:futures-margin-settlement]{futures margin and settlement}.

\qffield{Pitfalls} The legacy taxonomy listed venue risk (``outage, withdrawal halt, auto-deleveraging, or insolvency''), liquidation risk, stablecoin risk, bridge and settlement risk, and fragmentation (``basis and price gaps across venues when capital cannot move quickly''). Auto-deleveraging and liquidation are checked only for Deribit; outages and withdrawal halts were not checked; fragmentation is a market-structure fact (\hyperref[def:crypto-venue-instrument-structure]{crypto venue and instrument structure}), linked by Makarov and Schoar to capital controls rather than to slow crypto transfers. Stablecoin and bridge risks are stubbed for their own pathways.

\qffield{See also} \hyperref[mech:perpetual-funding]{perpetual funding}; \hyperref[def:crypto-venue-instrument-structure]{crypto venue and instrument structure}; \hyperref[mech:futures-margin-settlement]{futures margin and settlement}.

\qffield{Further reading} The legacy blocks it replaces, verbatim, are in \nolinkurl{derivations.tex} in the folder \nolinkurl{notes/artefacts/crypto-mm-microstructure/}.

\qffield{Index} crypto liquidation; maintenance margin; liquidation fee; insurance fund; bankrupt account; socialised loss; auto-deleveraging; ADL; liquidation waterfall; crypto custody; third-party custodian; stablecoin depeg; bridge risk; network settlement; venue insolvency; Deribit; VARA

\qffield{Status} Verification: partly checked 2026-10-15 (\texttt{gapclose-deribit-rules}, Rulebook version 2.0 applicable 13 August 2026, page updated 15 September 2026, retrieved 15 October 2026 through the help-centre API: Rules 1.1, 2 (guidance on custody), 10.1--10.11 and the definitions of Auto-Deleveraging and Insurance Fund; the Exchange Membership Terms not read; Rules 10.6--10.11 rechecked against the same version on 15 October 2026, correcting the socialised-loss trigger and its discretion). Stability: draft. Currency: Deribit as of October 2026.
\end{qfdefinition}

\begin{qfevidence}[Crypto carry on dated futures: size, drivers (small traders' demand for leverage) and limits to arbitrage (spot access, margin and cross-margining), 2019--2024]\label{ev:crypto-dated-futures-carry}
\qffield{Statement} Schmeling, Schrimpf and Todorov measure the annualised basis of one- and three-month constant-maturity bitcoin and ether futures over spot, using data from Binance, OKEx, FTX, Huobi, BitMEX, Deribit and the CME (the CME from August 2020) \citep[revision of 1 October 2025, abstract, sec.~1 and data section, PDF pp.~3--9]{roadmap-carry}.
\begin{itemize}
\item \emph{Size.} Carry ``can reach exceptionally high levels, sometimes exceeding 40\% per annum''. It averaged about 7 percent a year across exchanges from April 2019 to July 2024 (pp.~3 and~5).
\item \emph{Not fundamentals.} In $F_t=S_te^{(r_t+u_t-y_t)(T-t)}$, interest rates bear ``little to no relation'' to the carry and storage costs are negligible. That implies a large negative convenience yield: investors prefer being long futures to holding spot (p.~5).
\item \emph{Demand.} Higher carry goes with larger net long positions of smaller traders in the CFTC's Commitments of Traders data, who add positions in strong trends and periods of attention. Dealers and leveraged funds take the short side. Introducing CME micro bitcoin futures raised CME carry relative to other venues in a difference-in-differences test (pp.~6--7).
\item \emph{Limits to arbitrage.} Regulation long kept many professionals from holding spot, and margin frictions bind. With no cross-margining on traditional exchanges, a long spot position cannot collateralise a short futures position. The futures leg suffers severe drawdowns, and carry predicts liquidations of short futures: a 10 percent rise in standardised carry predicts short liquidations of 22 percent of open interest over the next month (p.~7). After spot bitcoin ETFs were introduced in 2024, carry fell by about 3 percentage points across exchanges and 5 more on the CME, 36 and 97 percent of the mean (p.~8).
\end{itemize}

\qffield{Conditions} The evidence is for fixed-date futures, which converge to spot at expiry; perpetual futures are discussed only through other papers (p.~8). The causal claims rest on two difference-in-differences designs around specific events. The exchange sample includes FTX, which failed in November 2022.

\qffield{Practice} A cash-and-carry trade that ``appears risk-free'' in the absence of frictions (p.~5) carries funding and liquidation risk before convergence. Size it against the margin path and the venue's liquidation rules, not against the annualised basis alone (\hyperref[def:crypto-trading-risk-layers]{crypto trading risk layers}). For perpetuals, the paper restates other work: He et al.\ find that perpetuals ``typically trade at a positive funding rate, that is, the long side pays the short side'', and Christin et al.\ that a carry trade short perpetual futures earns high Sharpe ratios (p.~8). These are not checked here.

\qffield{Pitfalls} The PDF was read from the Internet Archive's capture of 3 June 2026, because bis.org served a web page in place of the file on 2026-10-15 and again on 2026-10-17. The working paper's tables were not checked.

\qffield{See also} \hyperref[mech:perpetual-funding]{perpetual funding}; \hyperref[res:cost-of-carry-forward-price]{cost-of-carry forward price} (the pricing equation); \hyperref[def:contango-backwardation-stocks]{convenience yield in commodities} (opposite sign).

\qffield{Further reading} \texttt{roadmap-perpetual-fundamentals} for funding rates; \texttt{s11-r15}.

\qffield{Index} crypto carry; futures basis; cash-and-carry; convenience yield; inconvenience yield; limits to arbitrage; cross-margining; liquidations; micro bitcoin futures; spot bitcoin ETF; Commitments of Traders; leverage demand

\qffield{Status} Verification: partly checked 2026-10-17 (against \texttt{roadmap-carry}, BIS Working Paper 1087, revision of 1 October 2025, Internet Archive capture, PDF pp.~1--9, text layer). Stability: draft. Currency: evidence to July 2024.
\end{qfevidence}


\begin{qfdefinition}[Automated market maker pools: the constant-product rule and its fee, price impact and the cost of providing liquidity, and concentrated liquidity in Uniswap v3]\label{def:amm-pool-mechanics}
\qffield{Statement} A Uniswap pair pools reserves $x$ and $y$ of two tokens. Uniswap v1 and v2 maintain ``the invariant that the product of the reserves cannot decrease''; ``traders pay a 30-basis-point fee on trades, which goes to liquidity providers'' \citep[sec.~1, p.~1]{uniswap-v2-2020}. Version~2 enforces the fee through $(x_1-0.003\,x_{in})(y_1-0.003\,y_{in})\ge x_0y_0$, with $x_0,y_0$ the balances before and $x_1,y_1$ after the swap (sec.~3.2.1, eq.~(10), p.~7). Uniswap v3 keeps $x\cdot y=k$ but lets a liquidity provider confine liquidity $L=\sqrt k$ to a price range $[p_a,p_b]$, a \emph{position} \citep[sec.~2, eqs.~(2.1)--(2.2), p.~2, read from the page image]{s6-uniswap}. Its real reserves satisfy
\[ \Big(x+\frac{L}{\sqrt{p_b}}\Big)\big(y+L\sqrt{p_a}\big)=L^2, \]
so it trades like a constant-product pool with larger ``virtual reserves'' inside the range. Outside the range it holds one asset only, ``is no longer active, and no longer earns fees''. Fee tiers are set per pool, initially 0.05, 0.30 and 1 percent (sec.~1, p.~1).

\qffield{Notation} Derived (own; checked below), with price $p=y/x$ and fee $f$, $\gamma=1-f$: selling $\Delta x$ returns $\Delta y=y\gamma\Delta x/(x+\gamma\Delta x)$. A position's real reserves at price $p$ inside its range are $x=L(1/\sqrt p-1/\sqrt{p_b})$ and $y=L(\sqrt p-\sqrt{p_a})$.

\qffield{Conditions} The formulas are the protocols' design as described in their 2020 and 2021 whitepapers. Deployed pools, fee switches, later versions (UniswapX, v4) and other AMM designs differ. Gas costs, transaction ordering and the maximal extractable value they enable are outside these papers and were not read.

\qffield{Practice} Own checks, \nolinkurl{notes/artefacts/crypto-perpetuals/batch60-amm-checks.py}:
\begin{itemize}
\item \emph{Price impact.} Selling 10 ETH into a 1{,}000 ETH and 2{,}000{,}000 USDC pool returns 19{,}743.16 USDC, an execution price of 1{,}974.32 against a pool price of 2{,}000 before and 1{,}960.65 after. Before the fee the shortfall is $\gamma\Delta x/(x+\gamma\Delta x)$, 0.99 percent for a trade of 1 percent of the reserve. The product $k$ rises by the fee retained.
\item \emph{Cost of providing liquidity.} With arbitrageurs moving the pool to an external price $p_0r$, a full-range position is worth $2\sqrt r/(1+r)$ times the buy-and-hold value before fees: $-5.7$ percent for $r=2$ or $1/2$, and $-20$ percent for $r=4$. That is the ``impermanent loss'' that fee income must cover.
\item \emph{Concentration.} The same $L$ on $[1600,2500]$ at a price of 2{,}000 needs 94{,}427 of capital against 894{,}427 on the full range, 9.5 times less. Below $p_a$ the position is all X (50 units); above $p_b$, all Y (100{,}000).
\end{itemize}

\qffield{Pitfalls} The whitepapers are the developers' design descriptions, not audits of deployed code. The impermanent-loss expression assumes the pool is arbitraged to an outside price and ignores fees and gas; both are own derivations, not statements of the sources.

\qffield{See also} \hyperref[def:crypto-venue-instrument-structure]{crypto venues and instruments}; \hyperref[def:crypto-trading-risk-layers]{crypto trading risk layers}; \hyperref[ev:crypto-dated-futures-carry]{crypto carry}.

\qffield{Further reading} \texttt{roadmap-uniswapx}; \texttt{roadmap-cfmm-routing-efficient}.

\qffield{Index} automated market maker; AMM; constant product; x*y=k; Uniswap; liquidity pool; liquidity provider; swap fee; price impact; impermanent loss; concentrated liquidity; virtual reserves; range order

\qffield{Status} Verification: partly checked 2026-10-17 (against \texttt{uniswap-v2-2020}, the v2 Core whitepaper of March 2020, secs.~1 and~3.2.1, text layer; \texttt{s6-uniswap}, the v3 Core whitepaper of March 2021, secs.~1--2, text layer and a page image of p.~2; own numerical checks). Stability: draft. Currency: protocol designs of 2020--2021.
\end{qfdefinition}
